% RecSys Odyssey repository-backed lecture note.
% Add figures with: \coursefigure{figures/name.svg}{Caption}
\section{Retrieve for recall, rank for precision}

\begin{abstract}
Use different models for different scales instead of forcing one stage to do every job.
\end{abstract}

\subsection{Concept}
Blend fast candidate sources to preserve relevant and fresh items, then apply richer ranking features to the merged shortlist. Re-ranking handles diversity, constraints and slate interactions. Measure each boundary independently: candidate recall before ranking, ranking quality before selectors, and final outcomes after serving. Stage traces make it possible to explain where a valuable item was lost.

\begin{equation*}
catalogue \rightarrow multi-source retrieval \rightarrow ranker \rightarrow re-ranker/selectors \rightarrow response
\end{equation*}

\subsection{Key terms}
\begin{itemize}
\item \textbf{Stage contract.} The input, output, metric and latency responsibility of one pipeline stage.
\item \textbf{Candidate recall.} How many relevant items survive retrieval into the ranker.
\item \textbf{Stage trace.} A record of items entering, leaving and being dropped at every stage.
\end{itemize}
